% =============================================================================
%  lec24.tex —— 第 24 讲 模 II（底本 PDF 第 180–185 页）
% =============================================================================
\lecture{24}{模 II}
\begin{center}\small\color{black!55}（底本 PDF 第 180–185 页）\end{center}

本讲研究 $\infty$-操作子 $\OO^{\otimes}$ 上的 $\OO$-模。本讲始终假定 $\OO^{\otimes}$ 是连贯的 $\OO$-操作子，且 $q\colon\cC^{\otimes}\to\OO^{\otimes}$ 是 $\infty$-操作子映射。

\section{相对张量积}

\nitem\label{L24-1.1} 上一讲我们定义了一个 $\infty$-范畴 $\Mod^{\OO}(\cC)^{\otimes}$，它位于下述交换图中：
\begin{equation}\label{L24-eq1.1}
\begin{tikzcd}[column sep=4.8em,row sep=2.8em]
\Mod^{\OO}(\cC)^{\otimes} \arrow[r] \arrow[d] & \cC^{\otimes} \arrow[d, "q"] \\
\Alg_{/\OO}(\cC)\times\OO^{\otimes} \arrow[r] & \OO^{\otimes}.
\end{tikzcd}
\end{equation}

$\Mod^{\OO}(\cC)^{\otimes}$ 中的对象应视为三元组 $(A,X,M)$，其中
\begin{itemize}[leftmargin=2.2em,itemsep=0.1em]
\item $A\in\Alg_{/\OO}(\cC)$ 是 $\cC$ 中的一个 $\OO$-代数，
\item $X\in\OO^{\otimes}$ 是 $\OO^{\otimes}$ 中的一个\textbf{多色}，
\item $M$ 是 $\cC^{\otimes}_{X}$ 中的对象，带有型为 $X$ 的 $\OO$-$A$-模结构。
\end{itemize}
\eqref{L24-eq1.1} 中上方的水平函子把 $(A,X,M)$ 送到其底层对象 $M$，而左侧的竖直函子把该三元组送到 $(A,X)$。

\begin{theorem}\label{L24-1.2}（HA.3.3.3.9, 3.3.3.10）
设 $A\in\Alg_{/\OO}(\cC)$ 是 $\cC$ 中的 $\OO$-代数。则 $\Mod^{\OO}_{A}(\cC)^{\otimes}$ 是 $\infty$-操作子，且遗忘函子
\[
\Mod^{\OO}_{A}(\cC)^{\otimes}\to\cC^{\otimes}
\]
是 $\infty$-操作子映射。
\end{theorem}

\begin{proposition}\label{L24-1.3}（HA.3.4.2.1）
设 $q\colon\cC^{\otimes}\to\OO^{\otimes}$ 把 $\cC^{\otimes}$ 展示为 $\OO$-幺半 $\infty$-范畴。设 $A\defeq\mathbf{1}_{\cC}\in\Alg_{/\OO}(\cC)$ 是 $\cC$ 中的单位代数。则遗忘函子
\[
\Mod^{\OO}_{\mathbf{1}_{\cC}}(\cC)^{\otimes}\to\cC^{\otimes}
\]
是等价。
\end{proposition}

\begin{theorem}\label{L24-1.4}（HA.3.4.4.2）
设 $A\in\Alg_{/\OO}(\cC)$ 是 $\cC$ 中的 $\OO$-代数。假设：
\begin{enumerate}[label=(\roman*),leftmargin=2.4em,itemsep=0.1em]
\item $\infty$-操作子 $\OO^{\otimes}$ 是本质小的；
\item 函子 $q\colon\cC^{\otimes}\to\OO^{\otimes}$ 把 $\cC^{\otimes}$ 展示为可表现 $\OO$-幺半 $\infty$-范畴。
\end{enumerate}
则函子
\[
\Mod^{\OO}_{A}(\cC)^{\otimes}\to\OO^{\otimes}
\]
把 $\Mod^{\OO}_{A}(\cC)^{\otimes}$ 展示为可表现 $\OO$-幺半 $\infty$-范畴。
\end{theorem}

\begin{remark}\label{L24-1.5}
粗略地说，上述定理断言：$\OO$-幺半 $\infty$-范畴中 $A$-模的相对张量积，只要绝对张量积在余极限下表现良好，就存在。此外，这些相对张量积在 $A$-模的 $\infty$-范畴上定义了一个 $\OO$-幺半结构。
\end{remark}

\begin{remark}\label{L24-1.6}
当 $\OO^{\otimes}=\mathrm{Assoc}^{\otimes}$、$\mathrm{Comm}^{\otimes}$，或更一般地 $E_{k}^{\otimes}$（$k\ge1$）时，两个对象 $M,N\in\Mod^{\OO}_{A}(\cC)\defeq\Mod^{\OO}_{A}(\cC)^{\otimes}_{\langle 1\rangle}$ 的相对张量积 $M\otimes_{A}N$ 可以显式地计算为几何实现
\[
M\otimes N\longleftarrow M\otimes A\otimes N\longleftarrow M\otimes A\otimes A\otimes N\ \cdots
\]
称为 \textbf{Bar 构造}。更多细节见 HA.4.4 与 5.1.3。
\end{remark}

\begin{proposition}\label{L24-1.7}（HA.3.4.4.4）
在定理~\ref{L24-1.4} 的设定下，对任意多色 $X\in\OO^{\otimes}$，$\Mod^{\OO}_{A}(\cC)^{\otimes}$ 中的 $X$-单位对象的底层对象在 $\cC^{\otimes}$ 中典范同构于 $A(X)$。
\end{proposition}

\begin{remark}\label{L24-1.8}
更确切地说，$\Mod^{\OO}_{A}(\cC)^{\otimes}$ 中位于 $\langle 0\rangle\to X$ 上方的态射 $M_{0}\to M$ 把 $M$ 展示为 $X$-单位对象，当且仅当相应的函子
\[
\Delta^{1}\times_{\OO^{\otimes}}K_{\OO}\to\cC^{\otimes}
\]
把 $(\langle 0\rangle\to X)\in K_{\langle 0\rangle}$ 与 $(X\xrightarrow{\ \id_{X}\ }X)\in K_{X}$ 之间的态射送到 $\cC^{\otimes}_{X}$ 中的同构。
\end{remark}

\nitem\label{L24-1.9} 下述结果推广了交换代数中相应的著名结果。

\begin{proposition}\label{L24-1.10}（HA.3.4.1.7）
设 $A\in\Alg_{/\OO}(\cC)$ 是 $\cC$ 中的 $\OO$-代数。则存在与到 $\Alg_{/\OO}(\cC)$ 的遗忘函子相容的典范等价
\[
\Alg_{/\OO}(\Mod^{\OO}_{A}(\cC))\to\Alg_{/\OO}(\cC)_{A/}.
\]
\end{proposition}

\begin{remark}\label{L24-1.11}
粗略地说，上述等价由沿函子
\[
\OO^{\otimes}\times\Delta^{1}\to K_{\OO}
\]
的限制诱导，该函子把 $(X,0)$ 送到零态射 $\langle 0\rangle\to X$，把 $(X,1)$ 送到恒等态射 $X\xrightarrow{\ \id_{X}\ }X$。
\end{remark}

\nitem\label{L24-1.12} 我们也可以描述 $\Mod^{\OO}_{A}(\cC)^{\otimes}$ 中的模。

\begin{proposition}\label{L24-1.13}（HA.3.4.1.8）
设 $A\in\Alg_{/\OO}(\cC)$ 是 $\cC$ 中的 $\OO$-代数。则存在 $\infty$-范畴的典范拉回图：
\begin{center}
\begin{tikzcd}[column sep=2.4em,row sep=2.8em]
\Mod^{\OO}(\Mod^{\OO}_{A}(\cC))^{\otimes} \arrow[r] \arrow[d] & \Alg_{/\OO}(\Mod^{\OO}_{A}(\cC)) \arrow[r, "\simeq"] & \Alg_{/\OO}(\cC)_{A/} \arrow[d] \\
\Mod^{\OO}(\cC)^{\otimes} \arrow[rr] & & \Alg_{/\OO}(\cC).
\end{tikzcd}
\end{center}
\end{proposition}

\nitem\label{L24-1.14} 上述命题的含义可以通过它的下述推论来说明。

\begin{corollary}\label{L24-1.15}
设 $\phi\colon A\to B$ 是 $\Alg_{/\OO}(\cC)$ 中的态射。则存在典范等价
\[
\Mod^{\OO}_{B}(\Mod^{\OO}_{A}(\cC))^{\otimes}\xrightarrow{\ \simeq\ }\Mod^{\OO}_{B}(\cC)^{\otimes},
\]
其中在源中我们把 $B$ 视为 $\Mod^{\OO}_{A}(\cC)^{\otimes}$ 中的 $\OO$-代数。
\end{corollary}

\section{限制与极限}

\begin{theorem}\label{L24-2.1}（HA.3.4.3.4）
函子
\begin{equation}\label{L24-eq2.1}
\Mod^{\OO}(\cC)^{\otimes}\to\Alg_{/\OO}(\cC)
\end{equation}
是笛卡尔纤维化，且 $\Mod^{\OO}(\cC)^{\otimes}$ 中的态射是笛卡尔的，当且仅当其像在 $\cC^{\otimes}$ 中同构。
\end{theorem}

\begin{exercise}\label{L24-2.2}
设 $A\in\Alg_{/\OO}(\cC)$。证明遗忘函子 $\Mod^{\OO}_{A}(\cC)^{\otimes}\to\cC^{\otimes}$ 是保守的。
\end{exercise}

\begin{exercise}\label{L24-2.3}
证明 \eqref{L24-eq2.1} 的拉直
\[
\Alg_{/\OO}(\cC)^{\op}\to\Cat_{\infty}
\]
典范地穿过 $(\Cat_{\infty})_{/\cC^{\otimes}}$。
\end{exercise}

\begin{exercise}\label{L24-2.4}
设 $A\in\Alg_{/\OO}(\cC)$。证明 $\Mod^{\OO}_{A}(\cC)^{\otimes}$ 中的态射是惰性态射，当且仅当其像在 $\cC^{\otimes}$ 中是惰性态射。由此推出：对 $\Alg_{/\OO}(\cC)$ 中的态射 $\phi\colon A\to B$，反变传输函子
\[
\res_{\phi}\colon\Mod^{\OO}_{B}(\cC)^{\otimes}\to\Mod^{\OO}_{A}(\cC)^{\otimes}
\]
是 $\infty$-操作子映射。
\end{exercise}

\begin{remark}\label{L24-2.5}
粗略地说，上述结果说明：沿 $\OO$-代数之间的同态存在限制函子，而这些限制函子在相对张量积存在时，与相对张量积松弛地相容。
\end{remark}

\begin{exercise}\label{L24-2.6}
$\res_{\phi}$ 与等价
\[
\Mod^{\OO}_{B}(\Mod^{\OO}_{A}(\cC))^{\otimes}\xrightarrow{\ \simeq\ }\Mod^{\OO}_{B}(\cC)^{\otimes}
\]
之间是什么关系？
\end{exercise}

\nitem\label{L24-2.7} 下一个结果关于模的极限。应把它与第 22 讲定理 5.3 相比较，后者关于代数的极限。

\begin{theorem}\label{L24-2.8}（HA.3.4.3.6）
设 $A\in\Alg_{/\OO}(\cC)$ 是 $\cC$ 中的 $\OO$-代数。设 $X\in\OO^{\otimes}$ 是多色，$K$ 是单纯集，并假设：
\begin{enumerate}[label=(\roman*),leftmargin=2.4em,itemsep=0.1em]
\item 函子 $q\colon\cC^{\otimes}\to\OO^{\otimes}$ 是余笛卡尔纤维化；
\item 纤维 $\cC^{\otimes}_{X}$ 有 $K$-指标的极限。
\end{enumerate}
则 $\infty$-范畴 $\Mod^{\OO}_{A}(\cC)^{\otimes}_{X}$ 有 $K$-指标的极限，且遗忘函子
\[
\Mod^{\OO}_{A}(\cC)^{\otimes}_{X}\to\cC^{\otimes}_{X}
\]
保持并检测 $K$-指标的极限。
\end{theorem}

\begin{exercise}\label{L24-2.9}
在上述设定下，证明：对 $\Alg_{/\OO}(\cC)$ 中的任意态射 $\phi\colon A\to B$，函子
\[
\res_{\phi}\colon\Mod^{\OO}_{B}(\cC)^{\otimes}_{X}\to\Mod^{\OO}_{A}(\cC)^{\otimes}_{X}
\]
保持 $K$-指标的极限。
\end{exercise}

\begin{remark}\label{L24-2.10}
事实上，定理~\ref{L24-2.1} 与定理~\ref{L24-2.8} 可以合并为关于 $\Mod^{\OO}(\cC)^{\otimes}$ 中相对于投影
\[
\Mod^{\OO}(\cC)^{\otimes}\to\Alg_{/\OO}(\cC)\times\OO^{\otimes}
\]
的极限的单一陈述。更多细节见 HA.3.4.3.1。
\end{remark}

\section{诱导与余极限}

\nitem\label{L24-3.1} 练习~\ref{L24-2.9} 与伴随函子定理蕴含下述结果。

\begin{corollary}\label{L24-3.2}
假设：
\begin{enumerate}[label=(\roman*),leftmargin=2.4em,itemsep=0.1em]
\item $\infty$-操作子 $\OO^{\otimes}$ 是本质小的；
\item 函子 $q\colon\cC^{\otimes}\to\OO^{\otimes}$ 把 $\cC^{\otimes}$ 展示为可表现 $\OO$-幺半 $\infty$-范畴。
\end{enumerate}
则对任意多色 $X\in\OO^{\otimes}$，函子
\[
\Mod^{\OO}(\cC)^{\otimes}_{X}\to\Alg_{/\OO}(\cC)
\]
是余笛卡尔纤维化。
\end{corollary}

\begin{remark}\label{L24-3.3}
对 $\Alg_{/\OO}(\cC)$ 中的态射 $\phi\colon A\to B$，协变传输函子
\[
\ind_{\phi}\colon\Mod^{\OO}_{A}(\cC)^{\otimes}_{X}\to\Mod^{\OO}_{B}(\cC)^{\otimes}_{X}
\]
称为\textbf{诱导函子}。
\end{remark}

\nitem\label{L24-3.4} 下述结果说明，当 $\OO^{\otimes}=\mathrm{Comm}^{\otimes}$ 时，诱导函子是对称幺半的。

\begin{theorem}\label{L24-3.5}（HA.4.5.1.3）
设 $\OO^{\otimes}=\mathrm{Comm}^{\otimes}$。假设：
\begin{itemize}[leftmargin=2.2em,itemsep=0.1em]
\item 函子 $q\colon\cC^{\otimes}\to\mathrm{Comm}^{\otimes}$ 把 $\cC^{\otimes}$ 展示为与几何实现相容的对称幺半 $\infty$-范畴。
\end{itemize}
则函子
\[
\Mod^{\mathrm{Comm}}(\cC)^{\otimes}\to\Alg_{\mathrm{Comm}}(\cC)\times\OO^{\otimes}
\]
是余笛卡尔纤维化。
\end{theorem}

\begin{remark}\label{L24-3.6}
诱导函子
\[
\ind_{\phi}\colon\Mod^{\mathrm{Comm}}_{A}(\cC)\to\Mod^{\mathrm{Comm}}_{B}(\cC)
\]
把 $M$ 送到 $B\otimes_{A}M$，后者可以显式地计算为几何实现
\[
B\otimes M\longleftarrow B\otimes A\otimes N\longleftarrow B\otimes A\otimes A\otimes M\ \cdots.
\]
更多细节见 HA.4.5.3。
\end{remark}

\yizhu{注~\ref{L24-3.6} 中几何实现的第二项印作 $B\otimes A\otimes N$，其中 $N$ 应为 $M$：对照注~\ref{L24-1.6} 中相对张量积的 Bar 构造即知。译本照原文抄录符号，未作改动。\textbf{\mbox{[来源：I]}}}

\nitem\label{L24-3.7} 回忆：代数的一般余极限未必被遗忘函子保持。然而，对模而言这不会发生。

\begin{theorem}\label{L24-3.8}（HA.3.4.4.6）
设 $A\in\Alg_{/\OO}(\cC)$ 是 $\cC$ 中的 $\OO$-代数。设 $X\in\OO^{\otimes}$ 是多色，$K$ 是单纯集，并假设：
\begin{itemize}[leftmargin=2.2em,itemsep=0.1em]
\item 函子 $q\colon\cC^{\otimes}\to\OO^{\otimes}$ 把 $\cC^{\otimes}$ 展示为与 $K$-指标的余极限相容的 $\OO$-幺半 $\infty$-范畴\footnote{HA.3.4.4.6 是在假设 $q$ 与任意小余极限相容之下陈述的。然而，我们更弱的假设已足以得到下面的断言，即上引文献中的断言 (2)。}
\end{itemize}
则 $\infty$-范畴 $\Mod^{\OO}_{A}(\cC)^{\otimes}_{X}$ 有 $K$-指标的余极限，且遗忘函子
\[
\Mod^{\OO}_{A}(\cC)^{\otimes}_{X}\to\cC^{\otimes}_{X}
\]
保持并检测 $K$-指标的余极限。
\end{theorem}

\begin{remark}\label{L24-3.9}
请用类似于第 22 讲 §5.5 的论证，解释上述定理为何合理。
\end{remark}

\begin{remark}\label{L24-3.10}
事实上，定理~\ref{L24-1.4} 与定理~\ref{L24-3.8} 可以合并为关于 $\Mod^{\OO}_{A}(\cC)^{\otimes}$ 中相对于投影
\[
\Mod^{\OO}_{A}(\cC)^{\otimes}\to\OO^{\otimes}
\]
的余极限的单一陈述。更多细节见 HA.3.4.4.3。
\end{remark}

\section{左模与双模}

\nitem\label{L24-4.1} 下述结果说明，交换代数的 $\mathrm{Comm}$-模等同于其底结合代数的左模（或右模）。

\begin{proposition}\label{L24-4.2}（HA.4.5.1.4）
设 $\OO^{\otimes}=\mathrm{Comm}^{\otimes}$。存在 $\infty$-范畴的典范拉回图：
\begin{center}
\begin{tikzcd}[column sep=3.8em,row sep=2.8em]
\Mod^{\mathrm{Comm}}(\cC) \arrow[r] \arrow[d] & \Alg_{\mathrm{LMod}}(\cC) \arrow[d] \\
\Alg_{\mathrm{Comm}}(\cC) \arrow[r] & \Alg_{\mathrm{Assoc}}(\cC).
\end{tikzcd}
\end{center}
\end{proposition}

\nitem\label{L24-4.3} 下述结果说明，结合代数的 $\mathrm{Assoc}$-模等同于它的双模。

\begin{proposition}\label{L24-4.4}（HA.4.4.1.28）
设 $\OO^{\otimes}=\mathrm{Assoc}^{\otimes}$。存在 $\infty$-范畴的典范拉回图：
\begin{center}
\begin{tikzcd}[column sep=2.2em,row sep=2.8em]
\Mod^{\mathrm{Assoc}}(\cC) \arrow[r] \arrow[d] & \Alg_{\mathrm{BMod}/\mathrm{Assoc}}(\cC) \arrow[d] \\
\Alg_{/\mathrm{Assoc}}(\cC) \arrow[r] & \Alg_{/\mathrm{Assoc}}(\cC)\times\Alg_{/\mathrm{Assoc}}(\cC),
\end{tikzcd}
\end{center}
其中底部的水平函子是对角函子。
\end{proposition}

\appendixsection{再论双模}

\nitem\label{L24-A.1} 知道一个 $(A,B)$-双模，等价于知道一个右 $B$-模，连同与 $B$-作用交换的左 $A$-作用。用公式表示，
\[
{}_{A}\mathrm{BMod}_{B}(\cC)\simeq\mathrm{LMod}_{A}(\mathrm{RMod}_{B}(\cC)).
\]

\nitem\label{L24-A.2} $\Mod^{\mathrm{Assoc}}_{A}(\cC)$ 上的幺半结构可以推广为运算
\[
{}_{A_{0}}\mathrm{BMod}_{A_{1}}(\cC)\times{}_{A_{1}}\mathrm{BMod}_{A_{2}}(\cC)\times\cdots\times{}_{A_{n-1}}\mathrm{BMod}_{A_{n}}(\cC)\to{}_{A_{0}}\mathrm{BMod}_{A_{n}}(\cC),
\]
它们满足某些相容性。

\nitem\label{L24-A.3} \textbf{建议阅读}\quad HA.4.3。
